\section{CONFIG-010: Model Constraints Inconsistent Naming}
\label{sec:config-010}

\subsection*{Metadata}
\begin{tabular}{ll}
\textbf{Issue ID} & CONFIG-010 \\
\textbf{Severity} & Low \\
\textbf{Category} & Configuration \\
\textbf{Branch} & \branchref{fix/CONFIG-010-model-constraints-naming} \\
\textbf{Changes} & \branchdiff{fix/CONFIG-010-model-constraints-naming} \\
\end{tabular}

\subsection*{Executive Summary}
In \texttt{src/sdk.ts:286-292}, model constraints from SDK options are checked using both snake\_case and camelCase variants: \texttt{c.maxTokens} and \texttt{c.max\_tokens}, \texttt{c.topP} and \texttt{c.top\_p}, \texttt{c.topK} and \texttt{c.top\_k}. However, the \texttt{AgentManifest} interface in \texttt{src/loader.ts} only defines snake\_case variants. The dual-check pattern creates confusion about which convention is canonical and loses type safety through the \texttt{as any} cast.

The fix introduces a normalization function that maps both naming conventions to a canonical snake\_case form, removing the redundant dual checks and restoring type clarity.

\subsection*{Root Cause Analysis}

The original code checked both variants with no canonical form:

\begin{lstlisting}[language=TypeScript]
// src/sdk.ts:286-292 — Before fix
const constraints = options.constraints ?? loaded.manifest.model.constraints;
if (constraints) {
    const c = constraints as any;
    if (c.temperature !== undefined) modelOptions.temperature = c.temperature;
    if (c.maxTokens !== undefined) modelOptions.maxTokens = c.maxTokens;
    if (c.max_tokens !== undefined) modelOptions.maxTokens = c.max_tokens;
    if (c.topP !== undefined) modelOptions.topP = c.topP;
    if (c.top_p !== undefined) modelOptions.topP = c.top_p;
    if (c.topK !== undefined) modelOptions.topK = c.topK;
    if (c.top_k !== undefined) modelOptions.topK = c.top_k;
}
\end{lstlisting}

The manifest interface only defines snake\_case:

\begin{lstlisting}[language=TypeScript]
// src/loader.ts:38-44 — Manifest interface
model: {
    preferred: string;
    fallback: string[];
    constraints?: {
        temperature?: number;
        max_tokens?: number;    // snake_case
        top_p?: number;         // snake_case
        top_k?: number;         // snake_case
        stop_sequences?: string[];
    };
};
\end{lstlisting}

Problems with the dual-check approach:
\begin{itemize}
    \item \textbf{No canonical convention}: Config files use \texttt{max\_tokens}, SDK accepts \texttt{maxTokens}
    \item \textbf{Silent precedence}: If both \texttt{maxTokens} and \texttt{max\_tokens} are provided, \texttt{max\_tokens} wins because it is checked second
    \item \textbf{Lost type safety}: The \texttt{as any} cast bypasses TypeScript checking
\end{itemize}

\subsection*{Fix Implementation}

The fix introduces a normalization function with a mapping table:

\begin{lstlisting}[language=TypeScript]
// src/sdk.ts — After fix
const CONSTRAINT_MAPPING: Record<string, string> = {
    maxTokens: "max_tokens",
    max_tokens: "max_tokens",
    topP: "top_p",
    top_p: "top_p",
    topK: "top_k",
    top_k: "top_k",
    temperature: "temperature",
    stop_sequences: "stop_sequences",
    stopSequences: "stop_sequences",
};

function normalizeConstraints(raw: Record<string, any>): Record<string, any> {
    const result: Record<string, any> = {};
    for (const [key, value] of Object.entries(raw)) {
        const canonical = CONSTRAINT_MAPPING[key];
        if (canonical) result[canonical] = value;
    }
    return result;
}
\end{lstlisting}

The usage site is now clean and only references the canonical form:

\begin{lstlisting}[language=TypeScript]
// src/sdk.ts — After fix: usage
const constraints = options.constraints ?? loaded.manifest.model.constraints;
if (constraints) {
    const c = normalizeConstraints(constraints as any);
    if (c.temperature !== undefined) modelOptions.temperature = c.temperature;
    if (c.max_tokens !== undefined) modelOptions.maxTokens = c.max_tokens;
    if (c.top_p !== undefined) modelOptions.topP = c.top_p;
    if (c.top_k !== undefined) modelOptions.topK = c.top_k;
}
\end{lstlisting}

\subsection*{Affected Files}
\begin{itemize}
    \item \srcfile{src/sdk.ts} --- added \texttt{normalizeConstraints} function, removed dual checks
    \item \srcfile{src/\_\_tests\_\_/config-010.test.ts} --- new normalization tests
\end{itemize}

\subsection*{Verification}
\begin{lstlisting}[language=TypeScript]
// config-010.test.ts
test("normalizeConstraints maps camelCase to snake_case", () => {
    const result = normalizeConstraints(
        { maxTokens: 100, topP: 0.9, topK: 50 });
    deepStrictEqual(result, { max_tokens: 100, top_p: 0.9, top_k: 50 });
});

test("normalizeConstraints passes snake_case through unchanged", () => {
    const result = normalizeConstraints(
        { max_tokens: 100, top_p: 0.9, top_k: 50 });
    deepStrictEqual(result, { max_tokens: 100, top_p: 0.9, top_k: 50 });
});

test("normalizeConstraints prefers snake_case when both forms present", () => {
    const result = normalizeConstraints(
        { maxTokens: 50, max_tokens: 100 });
    deepStrictEqual(result, { max_tokens: 100 });
});
\end{lstlisting}
